\section{Fondasi Kecerdasan Buatan dan Basis Data Relasional}

Logika matematika bukan hanya sekedar pisau untuk memenggal arus *If-Else* di piranti lunak. Pada level yang jauh lebih makroskopis, segenap pilar arsitektur penyimpanan data raksasa masa kini (seperti *Database Oracle* dan *PostgreSQL*) maupun otak komputasi canggih *Artificial Intelligence* murni dibangun memuja fondasi Logika Predikat Orde Satu (\textit{First-Order Logic}) \cite{mit_6042,rosen2019}.

\subsection{1. Pemanggilan Data Presisi dengan Aljabar Relasional (SQL)}

Miliar baris \textit{record} tabungan nasabah di Bank tidak digali manual memakai kacamata pembesar manusia. Mesin mensurvei gunung data itu ditenagai **Aljabar Relasional** yang pada hakikatnya merupakan turunan langsung dari operasi *Konjungsi (AND), Disjungsi (OR)*, dan Kuantifikasi Predikat dari Logika Matematika di abad ke-19.

\begin{contoh}[Eksplorasi Query Bahasa Basis Data SQL]
Sebuah \textit{Query} dititahkan memburu profil tertentu: "Tampilkan seluruh Mahasiswa Fakultas MIPA \textbf{YANG MANA} umurnya belum menyentuh kepala Tiga \textbf{DAN} \textbf{TIDAK} sedang memiliki jejak skorsing kriminal."

Sintaks SQL (\textit{Structured Query Language}) murni mengadopsi struktur Proposisi Biner di atas:
\begin{lstlisting}[language=SQL, basicstyle=\ttfamily\small, frame=single, caption={Analisis Klausa \texttt{WHERE} dengan ikatan Logika Boolean Raksasa}]
SELECT nama_lengkap, nim 
FROM Tabel_Mahasiswa_Nasional
WHERE (fakultas = 'MIPA') 
  AND (usia_tahun < 30) 
  AND NOT (status_skorsing = 'Kriminal');
\end{lstlisting}
Selubung tebal kurung dan kasta \texttt{AND, OR, NOT} mutlak mengawal presisi *Database* membuang jutaan kandidat profil palsu \cite{rosen2019}.
\end{contoh}

\subsection{2. Nadi Penalaran pada *Rules Engine* Kecerdasan Buatan (AI)}

Memasuki ranah kecerdasan buatan, kita berjumpa dengan cabang *Sistem Pakar (Expert Systems)* yang menyandang julukan *Knowledge-based System Reasoning Engine*. Sistem-sistem kolosal peraba penyakit di dunia kedokteran hari ini bekerja menelaah deduksi \textit{Modus Ponens} dari triliunan fakta acak logis menjadi benang kesimpulan penentu nasib nyawa \cite{mit_6042,saylor_cs202}.

\begin{contoh}[Mesin Inferensi Medis Menarik Tali Diagnostik Maut]

\textbf{Basis Aturan (\textit{Rule Base}) dari Jurnal WHO yang diajarkan pada Mesin AI:}\\
Aturan 1: \textbf{JIKA} $Pasien_{x}$ memiliki Suhu di Atas 39 Celcius \textbf{DAN} Bercak Kulit Keras $\rightarrow$ \textbf{MAKA} \textit{Diagnosis Cacar Akut} (Valid!)\\
Aturan 2: \textbf{JIKA} $Pasien_{x}$ Diagnosa Cacar Akut \textbf{DAN} usia $X < 5$ \textbf{MAKA} $\rightarrow$ \textit{Prioritas Karantina ICUs!}

\textbf{Injeksi Fakta Premis Baru:} \\
Anak A berjemur demam $40$ Celcius, dipenuhi Bercak Kulit nan gatal. Si mungil baru rayakan ulang tahun usia $4$ tahun.

\textbf{Aksi \textit{Reasoning Engine} Membunuh Penyamaran (Inferensi):} 
Mesin tak perlu sekolah dokter 8 tahun. Secara *Modus Ponens* ia menggilas deduksi berjejaring, menyulut lonceng peringatan raksasa: \textbf{TARIK ANAK A MENUJU KARANTINA ICU RAWAT INAP!} 
Kemampuan program menyatukan titik serak tak kasat mata menjadi rentetan akibat sebab ini lahir membanggakan peradaban \textit{Inference Logic} yang Aristoteles rintis 20 abad lalu! \cite{stanford_cs103}
\end{contoh}
